\documentclass[a4paper,12pt]{article}

\usepackage[a4paper,margin=2.4cm]{geometry}
\usepackage{iftex}
\ifPDFTeX
  \usepackage[T2A]{fontenc}
  \usepackage[utf8]{inputenc}
  \usepackage[russian]{babel}
\else
  \usepackage[russian]{babel}
  \usepackage{fontspec}
  \setmainfont{Times New Roman}
\fi

\usepackage{amsmath}
\usepackage{amssymb}
\usepackage{enumitem}
\usepackage[hidelinks]{hyperref}

\setlength{\parindent}{0pt}
\setlength{\parskip}{0.35em}
\setlist[itemize]{leftmargin=1.5em,itemsep=0.15em,topsep=0.15em}
\setlist[enumerate]{leftmargin=1.7em,itemsep=0.15em,topsep=0.15em}
\setcounter{tocdepth}{2}

\newcommand{\R}{\mathbb{R}}
\newcommand{\ds}{\displaystyle}

\begin{document}

\begin{center}
{\LARGE Краткий конспект по курсу}\\[0.2em]
{\LARGE «Уравнения математической физики»}\\[0.6em]
{\large Коллоквиум 2}
\end{center}

\vspace{0.4em}
Эта версия повторяет структуру полного файла \texttt{umf.tex}, но оставляет только определения, ключевые формулы, условия применимости методов и короткие шаблоны решений.

\tableofcontents
\newpage

\section{Уравнения в частных производных первого порядка}

\subsection{1(a). Одномерное уравнение переноса с постоянным коэффициентом. Характеристики}

\textbf{Модель:}
\[
u_t+a u_x=0,
\qquad a=\mathrm{const}.
\]

\textbf{Характеристики} определяются уравнением
\[
\frac{dx}{dt}=a.
\]
Интегрирование даёт
\[
x=at+C,
\qquad x-at=C.
\]

\textbf{Вдоль характеристики}
\[
\frac{d}{dt}u(x(t),t)=u_t+u_x\frac{dx}{dt}=u_t+a u_x.
\]
Значит, для однородного уравнения
\[
\frac{d}{dt}u(x(t),t)=0,
\]
то есть решение постоянно на прямых $x-at=\mathrm{const}$.

\textbf{Физический смысл:} профиль не меняет форму, а просто сдвигается со скоростью $a$.

\textbf{Для запоминания:} перенос = движение вдоль прямых, без сглаживания и без рождения новых максимумов.

\subsection{1(b). Общее решение однородного уравнения переноса}

Так как $u$ постоянно на характеристиках, решение зависит только от их номера:
\[
u(x,t)=\Phi(x-at),
\]
где $\Phi$ --- произвольная функция.

\textbf{Проверка:}
\[
u_t=-a\Phi'(x-at),
\qquad
u_x=\Phi'(x-at),
\]
поэтому
\[
u_t+a u_x=0.
\]

\textbf{Итог:} всякое решение уравнения переноса есть сдвиг одной функции вдоль оси $x$.

\subsection{1(c). Задача Коши для уравнения переноса на всей прямой}

Рассматриваем
\[
u_t+a u_x=0,
\qquad x\in\R,\ t>0,
\qquad u(x,0)=\varphi(x).
\]

Подставляя $t=0$ в общее решение $u(x,t)=\Phi(x-at)$, получаем
\[
\Phi(x)=\varphi(x).
\]
Следовательно,
\[
u(x,t)=\varphi(x-at).
\]

\textbf{Смысл:} значение в точке $(x,t)$ берётся из начальной точки
\[
\xi=x-at.
\]

\textbf{Для классического решения} достаточно требовать, например, $\varphi\in C^1(\R)$.

\subsection{1(d). Начально-краевая задача на полупрямой. Метод характеристик}

Пусть
\[
u_t+a u_x=0,
\qquad x>0,\ t>0,
\qquad a>0,
\]
\[
u(x,0)=\varphi(x),
\qquad
u(0,t)=g(t).
\]

Характеристика, проходящая через $(x,t)$, либо пересекает начальную прямую $t=0$, либо границу $x=0$.

\textbf{Если $x\ge at$,} то характеристика приходит из точки $(x-at,0)$, поэтому
\[
u(x,t)=\varphi(x-at).
\]

\textbf{Если $0<x<at$,} то характеристика приходит из точки
\[
\left(0,t-\frac{x}{a}\right),
\]
поэтому
\[
u(x,t)=g\left(t-\frac{x}{a}\right).
\]

\textbf{Итак,}
\[
u(x,t)=
\begin{cases}
\varphi(x-at), & x\ge at,\\[0.4em]
g\left(t-\dfrac{x}{a}\right), & 0<x<at.
\end{cases}
\]

\textbf{Важно:} при $a>0$ граничное условие задаётся на входящей границе. При $a<0$ точка $x=0$ становится выходящей, и условие там задавать нельзя.

\subsection{1(e). Условие согласования в угловой точке}

В точке $(0,0)$ начальные и граничные данные должны не противоречить друг другу.

Для задачи
\[
u_t+a u_x=f(x,t),
\qquad
u(x,0)=\varphi(x),
\qquad
u(0,t)=g(t)
\]
\textbf{необходимое условие непрерывности:}
\[
\varphi(0)=g(0).
\]

Если требуется решение класса $C^1$, то уравнение в точке $(0,0)$ даёт ещё и \textbf{условие на производные:}
\[
g'(0)+a\varphi'(0)=f(0,0).
\]
Для однородного случая:
\[
g'(0)+a\varphi'(0)=0.
\]

\textbf{Для запоминания:} на углу сначала должны совпасть значения, а для гладкости --- ещё и уравнение.

\subsection{1(f). Полная производная вдоль характеристики и её связь с уравнением переноса}

Если характеристика задаётся законом $x=x(t)$, то
\[
\frac{d}{dt}u(x(t),t)=u_t+u_x\,x'(t).
\]
При $x'(t)=a$ получаем
\[
\frac{d}{dt}u(x(t),t)=u_t+a u_x.
\]

Поэтому уравнение
\[
u_t+a u_x=f(x,t)
\]
вдоль характеристики превращается в ОДУ
\[
\frac{d}{dt}u(x(t),t)=f(x(t),t).
\]

После интегрирования:
\[
u(x,t)=\varphi(x-at)+\int_0^t f(x-a(t-\tau),\tau)\,d\tau.
\]

\textbf{Ключевая идея метода характеристик:} PDE заменяется системой ОДУ вдоль специальных кривых.

\subsection{1(g). Конвективная (субстанциональная, материальная) производная в гидродинамике}

Пусть среда движется со скоростью $\mathbf{v}(x,t)$. Тогда для поля $q(x,t)$ \textbf{материальная производная} определяется так:
\[
\frac{Dq}{Dt}=q_t+\mathbf{v}\cdot\nabla q.
\]

\textbf{Смысл:}
\begin{itemize}
  \item $q_t$ --- локальное изменение в фиксированной точке;
  \item $\mathbf{v}\cdot\nabla q$ --- изменение из-за переноса частицей среды.
\end{itemize}

Если
\[
\frac{Dq}{Dt}=0,
\]
то величина $q$ сохраняется вдоль траекторий частиц.

\textbf{Частный случай:} при одномерном движении $\mathbf{v}=(a)$ получаем оператор
\[
\frac{D}{Dt}=\partial_t+a\partial_x,
\]
то есть именно оператор уравнения переноса.

\section{Практические задачи к разделу 1}

\subsection{Задача 1(a). Уравнение переноса с коэффициентом 2}

Дано:
\[
u_t+2u_x=0.
\]

\textbf{Характеристики:}
\[
\frac{dx}{dt}=2,
\qquad
x-2t=C.
\]

\textbf{Вдоль характеристики:}
\[
\frac{d}{dt}u(x(t),t)=u_t+2u_x=0.
\]
Следовательно,
\[
u(x,t)=\Phi(x-2t).
\]

Если задача понимается как задача Коши с условием
\[
u(x,0)=\varphi(x),
\]
то
\[
u(x,t)=\varphi(x-2t).
\]

\textbf{Ответ:}
\[
u(x,t)=\Phi(x-2t),
\qquad
u(x,t)=\varphi(x-2t)\ \text{при заданном }u(x,0)=\varphi(x).
\]

\subsection{Задача 1(b). Полная производная вдоль характеристики}

Для уравнения
\[
u_t+a u_x=0,
\qquad
u(x,0)=\varphi(x)
\]
характеристики имеют вид
\[
x=at+C.
\]

По правилу цепочки:
\[
\frac{d}{dt}u(x(t),t)=u_t+u_x\frac{dx}{dt}=u_t+a u_x.
\]
По уравнению это равно нулю, значит
\[
\frac{du}{dt}=0.
\]

Следовательно,
\[
u(x,t)=\varphi(x-at).
\]

\textbf{Ответ:}
\[
\frac{du}{dt}=u_t+a u_x,
\qquad
\frac{du}{dt}=0
\]
вдоль характеристик решения уравнения переноса.

\section{Метод энергетических неравенств}

\subsection{2(a). Суть энергетического метода}

Идея метода: \textbf{умножить уравнение на само решение} (или на подходящую производную), проинтегрировать по пространству и получить тождество для нормы решения.

Для параболического уравнения типа
\[
u_t-a u_{xx}+\gamma u=f
\]
после умножения на $2u$ обычно возникает структура
\[
\frac{d}{dt}\|u\|^2+\text{неотрицательные слагаемые}=2(f,u).
\]

\textbf{Зачем это нужно:}
\begin{itemize}
  \item получить априорные оценки;
  \item доказать единственность;
  \item доказать устойчивость по данным;
  \item контролировать гладкость через нормы.
\end{itemize}

\subsection{2(b). Энергетическое тождество для \texorpdfstring{$u_t-a u_{xx}+\gamma u=f$}{ut-a uxx+gamma u=f}}

Пусть
\[
u_t-a u_{xx}+\gamma u=f,
\qquad a>0,\ \gamma\ge 0,
\]
на полупрямой $x>0$, причём
\[
u(0,t)=0.
\]

Умножаем уравнение на $2u$ и интегрируем по $x$:
\[
2\int_0^\infty u u_t\,dx
-2a\int_0^\infty u u_{xx}\,dx
+2\gamma\int_0^\infty u^2\,dx
=2\int_0^\infty f u\,dx.
\]

Здесь
\[
2\int_0^\infty u u_t\,dx=\frac{d}{dt}\|u(\cdot,t)\|^2,
\qquad
\|u(\cdot,t)\|^2=\int_0^\infty u^2(x,t)\,dx.
\]

После интегрирования по частям и учёта $u(0,t)=0$:
\[
-2a\int_0^\infty u u_{xx}\,dx=2a\|u_x(\cdot,t)\|^2.
\]

\textbf{Итоговое тождество:}
\[
\frac{d}{dt}\|u(\cdot,t)\|^2
+2a\|u_x(\cdot,t)\|^2
+2\gamma\|u(\cdot,t)\|^2
=2(f(\cdot,t),u(\cdot,t)).
\]

\subsection{2(c). Как из энергетического тождества получается априорная оценка}

Из тождества
\[
\frac{d}{dt}\|u\|^2+2a\|u_x\|^2+2\gamma\|u\|^2=2(f,u)
\]
отбрасываем неотрицательные слагаемые слева:
\[
\frac{d}{dt}\|u\|^2\le 2(f,u).
\]

По неравенству Коши--Буняковского
\[
|(f,u)|\le \|f\|\,\|u\|.
\]
Если обозначить
\[
y(t)=\|u(\cdot,t)\|,
\]
то получаем
\[
y'(t)\le \|f(\cdot,t)\|.
\]
Интегрирование от $0$ до $t$ даёт
\[
\|u(\cdot,t)\|
\le
\|\varphi\|
+\int_0^t \|f(\cdot,\tau)\|\,d\tau.
\]

\textbf{Смысл:} норма решения контролируется начальными данными и правой частью.

\subsection{2(d). Единственность и устойчивость из априорной оценки}

Пусть $u_1$ и $u_2$ --- два решения одной задачи. Для разности
\[
w=u_1-u_2
\]
получаем однородную задачу
\[
w_t-a w_{xx}+\gamma w=0,
\qquad
w(x,0)=0,
\qquad
w(0,t)=0.
\]

Энергетическое тождество принимает вид
\[
\frac{d}{dt}\|w\|^2+2a\|w_x\|^2+2\gamma\|w\|^2=0.
\]
Интегрирование по времени даёт
\[
\|w(\cdot,t)\|^2=0.
\]
Следовательно,
\[
w\equiv 0,
\qquad
u_1\equiv u_2.
\]

\textbf{Устойчивость} получается тем же способом: разность решений оценивается через разность начальных и граничных данных, а также правых частей.

\subsection{2(e). Многомерный случай}

В области $\Omega\subset\R^n$ рассматривают
\[
u_t-a\Delta u+\gamma u=f(x,t),
\qquad
u|_{\partial\Omega}=0.
\]

Умножение на $2u$ и формула Грина дают
\[
\frac{d}{dt}\|u(\cdot,t)\|_{L^2(\Omega)}^2
+2a\|\nabla u(\cdot,t)\|_{L^2(\Omega)}^2
+2\gamma\|u(\cdot,t)\|_{L^2(\Omega)}^2
=2(f(\cdot,t),u(\cdot,t))_{L^2(\Omega)}.
\]

\textbf{Разница с одномерным случаем только в том, что}
\[
u_x^2 \ \leadsto\ |\nabla u|^2.
\]
Дальше схема доказательства априорных оценок и единственности полностью та же.

\section{Практические задачи к разделу 2}

\subsection{Задача 2(a). Вывести априорную оценку}

Для задачи
\[
u_t-a u_{xx}+\gamma u=f,
\qquad
u(0,t)=0,
\qquad
u(x,0)=\varphi(x)
\]
умножаем уравнение на $2u$ и интегрируем по $x$. Получаем
\[
\frac{d}{dt}\|u\|^2+2a\|u_x\|^2+2\gamma\|u\|^2=2(f,u).
\]

Отбрасывая неотрицательные члены слева и применяя
\[
|(f,u)|\le \|f\|\,\|u\|,
\]
получаем
\[
\frac{d}{dt}\|u\|^2\le 2\|f\|\,\|u\|.
\]

Обозначив $y(t)=\|u(\cdot,t)\|$, имеем
\[
y'(t)\le \|f(\cdot,t)\|.
\]
Интегрирование даёт
\[
\|u(\cdot,t)\|
\le
\|\varphi\|
+\int_0^t \|f(\cdot,\tau)\|\,d\tau.
\]

\textbf{Ответ:}
\[
\|u(\cdot,t)\|
\le
\|\varphi\|
+\int_0^t \|f(\cdot,\tau)\|\,d\tau.
\]

\subsection{Задача 2(b). Доказать единственность с помощью энергетического тождества}

Пусть $u_1$ и $u_2$ --- два решения. Тогда для
\[
w=u_1-u_2
\]
получаем
\[
w_t-a w_{xx}+\gamma w=0,
\qquad
w(0,t)=0,
\qquad
w(x,0)=0.
\]

Энергетическое тождество:
\[
\frac{d}{dt}\|w\|^2+2a\|w_x\|^2+2\gamma\|w\|^2=0.
\]
После интегрирования по времени:
\[
\|w(\cdot,t)\|^2=0.
\]
Значит,
\[
w\equiv 0,
\qquad
u_1\equiv u_2.
\]

\textbf{Ответ:} решение единственно.

\section{Уравнение Бюргерса и градиентная катастрофа}

\subsection{3(a). Вязкое и невязкое уравнение Бюргерса. Отличие от линейного переноса}

\textbf{Линейный перенос:}
\[
u_t+a u_x=0.
\]
Скорость $a$ фиксирована, характеристики параллельны:
\[
x-at=\mathrm{const}.
\]

\textbf{Невязкий Бюргерс:}
\[
u_t+u u_x=0.
\]
Теперь скорость равна самому решению, поэтому разные части профиля движутся с разными скоростями.

\textbf{Вязкий Бюргерс:}
\[
u_t+u u_x=\nu u_{xx},
\qquad \nu>0.
\]
Здесь к нелинейному переносу добавляется диффузия.

\textbf{Главное отличие от линейного переноса:}
\begin{itemize}
  \item для линейного переноса характеристики не пересекаются;
  \item для невязкого Бюргерса характеристики могут пересекаться;
  \item вязкость $\nu u_{xx}$ сглаживает решение и мешает образованию разрыва.
\end{itemize}

\subsection{3(b). Градиентная катастрофа (опрокидывание волны)}

Для задачи Коши
\[
u_t+u u_x=0,
\qquad
u(x,0)=\varphi(x)
\]
характеристики задаются формулами
\[
u=\varphi(\xi),
\qquad
x=\xi+t\varphi(\xi).
\]

Дифференцируя по параметру $\xi$, получаем
\[
x_\xi=1+t\varphi'(\xi).
\]
Пока
\[
1+t\varphi'(\xi)\ne 0
\]
для всех $\xi$, отображение $\xi\mapsto x$ взаимно однозначно и решение гладкое.

Производная решения имеет вид
\[
u_x=\frac{\varphi'(\xi)}{1+t\varphi'(\xi)}.
\]
Если знаменатель обращается в нуль, то
\[
|u_x|\to\infty.
\]
Это и есть \textbf{градиентная катастрофа}.

Если
\[
\min \varphi'(\xi)<0,
\]
то первый момент катастрофы:
\[
t_*=-\frac{1}{\min \varphi'(\xi)}.
\]

\subsection{3(c). Характеристики для уравнения Бюргерса и их пересечение}

Для невязкого уравнения Бюргерса характеристическая система:
\[
\frac{dx}{dt}=u,
\qquad
\frac{du}{dt}=0.
\]
Отсюда
\[
u=\varphi(\xi),
\qquad
x=\xi+t\varphi(\xi).
\]

\textbf{Пересечение характеристик} означает, что двум различным $\xi_1\ne \xi_2$ соответствует одна и та же точка $(x,t)$. Аналитически это эквивалентно условию
\[
x_\xi=1+t\varphi'(\xi)=0.
\]

\textbf{До момента пересечения} решение классическое, \textbf{после него} нужно переходить к слабым решениям и ударным волнам.

\subsection{3(d). Как вязкость предотвращает градиентную катастрофу}

При $\nu>0$ уравнение
\[
u_t+u u_x=\nu u_{xx}
\]
имеет параболический сглаживающий член $\nu u_{xx}$.

\textbf{Качественный механизм:}
\begin{itemize}
  \item нелинейность стремится круто наклонить профиль;
  \item диффузия выравнивает резкие перепады;
  \item вместо разрыва возникает гладкий переходный слой конечной толщины.
\end{itemize}

\textbf{Полезная формула:} преобразование Хопфа--Коула
\[
u=-2\nu \frac{\theta_x}{\theta}
\]
сводит уравнение Бюргерса к уравнению теплопроводности
\[
\theta_t=\nu \theta_{xx}.
\]
Это ещё раз показывает, что вязкий Бюргерс наследует сглаживающий эффект параболического уравнения.

\subsection{3(e). Ударная волна и её связь с градиентной катастрофой}

Невязкий Бюргерс можно переписать в законе сохранения:
\[
u_t+\left(\frac{u^2}{2}\right)_x=0.
\]

После пересечения характеристик классическое решение перестаёт существовать, и возникает \textbf{ударная волна} --- разрыв слабого решения.

Скорость разрыва определяется условием Ранкина--Гюгонио:
\[
s=\frac{[u^2/2]}{[u]}=\frac{u_-+u_+}{2},
\]
где $u_-$ и $u_+$ --- левые и правые значения на фронте.

\textbf{Энтропийная допустимость для Бюргерса:}
\[
u_->u_+.
\]
Именно такой разрыв возникает как предел вязких решений при $\nu\to 0+$.

\section{Практические задачи к разделу 3}

\subsection{Задача 3(a). Невязкое уравнение Бюргерса при \texorpdfstring{$u(x,0)=\sin x$}{u(x,0)=sin x}}

Рассматриваем
\[
u_t+u u_x=0,
\qquad
u(x,0)=\sin x.
\]

Характеристическая система:
\[
\frac{dx}{dt}=u,
\qquad
\frac{du}{dt}=0.
\]
Если характеристика стартует из точки $\xi$, то
\[
u=\sin \xi,
\qquad
x=\xi+t\sin \xi.
\]

Следовательно,
\[
u(x,t)=\sin \xi,
\qquad
x=\xi+t\sin \xi.
\]
Это параметрическая запись решения. Неявно:
\[
u=\sin(x-tu).
\]

Момент катастрофы определяется из
\[
x_\xi=1+t\cos\xi.
\]
Первое обращение в нуль:
\[
1-t=0,
\qquad
t_*=1.
\]

\textbf{Ответ:}
\[
u(x,t)=\sin \xi,
\qquad
x=\xi+t\sin \xi,
\qquad
0\le t<1,
\]
или
\[
u=\sin(x-tu),
\qquad
t_*=1.
\]

\subsection{Задача 3(b). Почему при \texorpdfstring{$u_t+u u_x=\nu u_{xx}$}{ut+u ux=nu uxx} катастрофы нет}

При $\nu>0$ появляется диффузионный член
\[
\nu u_{xx}.
\]
Он сглаживает растущий градиент и не даёт фронту стать вертикальным за конечное время.

Поэтому:
\begin{itemize}
  \item характеристики невязкой модели перестают быть точным инструментом;
  \item вместо разрыва возникает гладкий переходный слой;
  \item решение ведёт себя как нелинейная диффузия.
\end{itemize}

\textbf{Короткий вывод:} катастрофу предотвращает вязкость.

\section{Метод Фурье для параболического уравнения (теплопроводности)}

\subsection{4(a). Первая краевая задача для одномерного уравнения теплопроводности}

На отрезке $0<x<l$ рассматривается задача
\[
u_t=a^2u_{xx}+f(x,t),
\qquad 0<t\le T,
\]
\[
u(0,t)=g_1(t),
\qquad
u(l,t)=g_2(t),
\qquad
u(x,0)=\varphi(x).
\]

\textbf{Это первая краевая задача,} потому что на границе задаются значения самой функции $u$.

\textbf{Для классического решения} нужны условия согласования:
\[
\varphi(0)=g_1(0),
\qquad
\varphi(l)=g_2(0).
\]

\textbf{Основной модельный случай для метода Фурье:}
\[
u_t=a^2u_{xx},
\qquad
u(0,t)=u(l,t)=0,
\qquad
u(x,0)=\varphi(x).
\]

\subsection{4(b). Как применяется метод разделения переменных и какая возникает спектральная задача}

После сведения к однородным граничным условиям ищем решение в виде
\[
u(x,t)=X(x)T(t).
\]

Подстановка в
\[
u_t=a^2u_{xx}
\]
даёт
\[
X T'=a^2 X'' T.
\]
Делим на $a^2XT$:
\[
\frac{T'}{a^2T}=\frac{X''}{X}=-\lambda.
\]

Получаем две ОДУ:
\[
T'+a^2\lambda T=0,
\qquad
X''+\lambda X=0.
\]
Из условий
\[
u(0,t)=u(l,t)=0
\]
следует спектральная задача
\[
X''+\lambda X=0,
\qquad
X(0)=X(l)=0.
\]

\subsection{4(c). Собственные значения и собственные функции при условиях Дирихле}

Для задачи
\[
X''+\lambda X=0,
\qquad
X(0)=X(l)=0
\]
нет нетривиальных решений при $\lambda\le 0$.

При
\[
\lambda=\mu^2>0
\]
получаем
\[
X(x)=A\cos \mu x+B\sin \mu x.
\]
Из $X(0)=0$ следует $A=0$, а из $X(l)=0$:
\[
\sin \mu l=0
\quad\Longrightarrow\quad
\mu l=k\pi,\quad k=1,2,\dots
\]

\textbf{Собственные значения и функции:}
\[
\lambda_k=\left(\frac{k\pi}{l}\right)^2,
\qquad
X_k(x)=\sin\frac{k\pi x}{l},
\qquad
k=1,2,\dots
\]

\subsection{4(d). Вид решения в форме ряда Фурье и формулы для коэффициентов}

Для однородной задачи Дирихле решение записывается как
\[
u(x,t)=\sum_{k=1}^{\infty} c_k
e^{-a^2\left(\frac{k\pi}{l}\right)^2 t}
\sin\frac{k\pi x}{l}.
\]

Коэффициенты определяются разложением начальной функции:
\[
\varphi(x)\sim \sum_{k=1}^{\infty} c_k \sin\frac{k\pi x}{l},
\]
\[
c_k=\frac{2}{l}\int_0^l \varphi(x)\sin\frac{k\pi x}{l}\,dx.
\]

\textbf{Логика метода:} сначала строим базис из собственных функций, потом проектируем начальные данные на этот базис.

\subsection{4(e). Почему для параболического уравнения нужно одно начальное условие, а для гиперболического --- два}

Уравнение теплопроводности имеет \textbf{первый порядок по времени:}
\[
u_t=a^2u_{xx}+\dots
\]
Поэтому достаточно одного условия:
\[
u(x,0)=\varphi(x).
\]

Волновое уравнение имеет \textbf{второй порядок по времени:}
\[
u_{tt}=a^2u_{xx}+\dots
\]
Поэтому нужно два начальных условия:
\[
u(x,0)=\varphi(x),
\qquad
u_t(x,0)=\psi(x).
\]

\textbf{Коротко:} число начальных условий равно порядку уравнения по времени.

\subsection{4(f). Поведение решения при \texorpdfstring{$t\to\infty$}{t->infinity}: затухание мод}

Каждая гармоника в решении имеет множитель
\[
e^{-a^2\left(\frac{k\pi}{l}\right)^2 t}.
\]
При $t\to\infty$ все эти множители стремятся к нулю, причём высокие моды затухают быстрее.

Если нет источника и заданы однородные граничные условия, то
\[
u(x,t)\to 0.
\]

\textbf{Долговременное поведение} определяется первой модой:
\[
\sin\frac{\pi x}{l}.
\]

\section{Практические задачи к разделу 4}

\subsection{Задача 4(a). Решить \texorpdfstring{$u_t=u_{xx}$}{ut=uxx} на \texorpdfstring{$(0,\pi)$}{(0,pi)} при \texorpdfstring{$u(0,t)=u(\pi,t)=0$}{u(0,t)=u(pi,t)=0}, \texorpdfstring{$u(x,0)=\sin 2x$}{u(x,0)=sin 2x}}

Общий вид решения:
\[
u(x,t)=\sum_{k=1}^{\infty} a_k e^{-k^2 t}\sin kx.
\]

Начальная функция уже является одной собственной модой:
\[
\sin 2x.
\]
Поэтому
\[
a_2=1,
\qquad
a_k=0\quad (k\ne 2).
\]

\textbf{Ответ:}
\[
u(x,t)=e^{-4t}\sin 2x.
\]

\subsection{Задача 4(b). Для \texorpdfstring{$u(x,0)=x(\pi-x)$}{u(x,0)=x(pi-x)} найти коэффициенты Фурье}

Ищем решение в виде
\[
u(x,t)=\sum_{k=1}^{\infty} b_k e^{-k^2 t}\sin kx,
\]
где
\[
b_k=\frac{2}{\pi}\int_0^\pi x(\pi-x)\sin kx\,dx.
\]

После интегрирования по частям получается
\[
b_k=\frac{4(1-(-1)^k)}{\pi k^3}
=
\begin{cases}
\dfrac{8}{\pi k^3}, & k \text{ нечётно},\\[0.5em]
0, & k \text{ чётно}.
\end{cases}
\]

\textbf{Ответ:}
\[
u(x,t)=\frac{8}{\pi}\sum_{m=0}^{\infty}
\frac{e^{-(2m+1)^2 t}}{(2m+1)^3}\sin\bigl((2m+1)x\bigr).
\]

\subsection{Задача 4(c). Почему при \texorpdfstring{$t>0$}{t>0} решение быстро становится гладким}

В ряде Фурье каждая мода умножается на экспоненту
\[
e^{-c k^2 t},
\qquad c>0.
\]
При любом фиксированном $t>0$ эта экспонента подавляет высокочастотные гармоники быстрее любой степени $k$.

Поэтому ряды для производных по $x$ и по $t$ тоже сходятся, и решение мгновенно становится очень гладким.

\textbf{Ответ:} уравнение теплопроводности обладает \textbf{сглаживающим эффектом}.

\section{Метод Фурье для внутренней задачи уравнения Лапласа}

\subsection{5(a). Задача Дирихле для уравнения Лапласа в круге}

Пусть
\[
\Omega_a=\{(x,y)\in\R^2:\ x^2+y^2<a^2\}.
\]
Нужно найти
\[
u\in C^2(\Omega_a)\cap C(\overline{\Omega_a}),
\qquad
\Delta u=0\ \text{в }\Omega_a,
\qquad
u=g\ \text{на }\partial\Omega_a.
\]

В полярных координатах граничное условие записывается как
\[
u(a,\varphi)=g(\varphi),
\qquad
0\le \varphi\le 2\pi.
\]

\textbf{Условие на данные:} функция $g$ должна быть $2\pi$-периодической.

\textbf{Единственность} следует из принципа максимума.

\subsection{5(b). Метод разделения переменных в полярных координатах. Радиальная и угловая части}

В полярных координатах
\[
\Delta u=u_{\rho\rho}+\frac{1}{\rho}u_\rho+\frac{1}{\rho^2}u_{\varphi\varphi}.
\]

Ищем решение в виде
\[
u(\rho,\varphi)=R(\rho)\Phi(\varphi).
\]
После подстановки и разделения переменных получаем
\[
\frac{\rho^2 R''+\rho R'}{R}=-\frac{\Phi''}{\Phi}=\lambda.
\]

Отсюда возникают уравнения
\[
\Phi''+\lambda\Phi=0,
\qquad
\rho^2R''+\rho R'-\lambda R=0.
\]

Для угловой части добавляется периодичность:
\[
\Phi(0)=\Phi(2\pi),
\qquad
\Phi'(0)=\Phi'(2\pi).
\]

\subsection{5(c). Собственные значения и собственные функции угловой части}

Спектральная задача
\[
\Phi''+\lambda\Phi=0,
\qquad
\Phi(0)=\Phi(2\pi),
\qquad
\Phi'(0)=\Phi'(2\pi)
\]
имеет собственные значения
\[
\lambda_k=k^2,
\qquad
k=0,1,2,\dots
\]

\textbf{Собственные функции:}
\[
\Phi_0(\varphi)=1,
\]
\[
\Phi_k^{(1)}(\varphi)=\cos k\varphi,
\qquad
\Phi_k^{(2)}(\varphi)=\sin k\varphi,
\qquad
k\ge 1.
\]

\textbf{Причина:} на окружности естественным базисом являются гармоники по углу.

\subsection{5(d). Общее решение в виде ряда по косинусам и синусам}

Если
\[
g(\varphi)\sim \frac{\alpha_0}{2}
+\sum_{k=1}^{\infty}\bigl(\alpha_k\cos k\varphi+\beta_k\sin k\varphi\bigr),
\]
то \textbf{регулярное в центре} решение внутренней задачи Дирихле имеет вид
\[
u(\rho,\varphi)=\frac{\alpha_0}{2}
+\sum_{k=1}^{\infty}
\left(\frac{\rho}{a}\right)^k
\bigl(\alpha_k\cos k\varphi+\beta_k\sin k\varphi\bigr).
\]

\textbf{Почему нет членов} $\rho^{-k}$ и $\ln\rho$: они сингулярны при $\rho=0$, а внутреннее решение должно оставаться ограниченным в центре.

\section{Практические задачи к разделу 5}

\subsection{Задача 5(a). Решить задачу Дирихле в круге радиуса \texorpdfstring{$R$}{R} при \texorpdfstring{$u(R,\varphi)=\cos\varphi$}{u(R,phi)=cos phi}}

Граничная функция уже является одной гармоникой:
\[
g(\varphi)=\cos\varphi.
\]
Значит,
\[
\alpha_1=1,
\qquad
\alpha_k=0\ (k\ne 1),
\qquad
\beta_k=0.
\]

В общем ряде остаётся только мода $k=1$:
\[
u(\rho,\varphi)=\left(\frac{\rho}{R}\right)\cos\varphi.
\]

\textbf{Ответ:}
\[
u(\rho,\varphi)=\frac{\rho}{R}\cos\varphi.
\]
В декартовых координатах:
\[
u(x,y)=\frac{x}{R}.
\]

\section{Ряд Фурье и разложение по специальному базису}

\subsection{6(a). Определение тригонометрического ряда Фурье и формулы коэффициентов}

На $[0,2\pi]$ тригонометрический ряд Фурье имеет вид
\[
g(\varphi)\sim \frac{\alpha_0}{2}
+\sum_{k=1}^{\infty}\bigl(\alpha_k\cos k\varphi+\beta_k\sin k\varphi\bigr).
\]

Коэффициенты:
\[
\alpha_0=\frac{1}{\pi}\int_0^{2\pi} g(\psi)\,d\psi,
\]
\[
\alpha_k=\frac{1}{\pi}\int_0^{2\pi} g(\psi)\cos k\psi\,d\psi,
\qquad
\beta_k=\frac{1}{\pi}\int_0^{2\pi} g(\psi)\sin k\psi\,d\psi.
\]

На отрезке $(0,l)$ часто используют специальные ряды:
\[
f(x)\sim \sum_{k=1}^{\infty}b_k\sin\frac{k\pi x}{l},
\qquad
b_k=\frac{2}{l}\int_0^l f(x)\sin\frac{k\pi x}{l}\,dx,
\]
\[
f(x)\sim \frac{a_0}{2}+\sum_{k=1}^{\infty}a_k\cos\frac{k\pi x}{l},
\qquad
a_k=\frac{2}{l}\int_0^l f(x)\cos\frac{k\pi x}{l}\,dx.
\]

\subsection{6(b). Полнота и ортогональность. Примеры ортогональных систем}

Система $\{\phi_k\}$ ортогональна, если
\[
\langle \phi_k,\phi_m\rangle=0
\qquad (k\ne m),
\]
а полна, если по ней можно разлагать произвольные функции из рассматриваемого пространства.

\textbf{Примеры:}
\begin{itemize}
  \item $\sin\dfrac{k\pi x}{l}$ на $(0,l)$ --- для условий Дирихле;
  \item $1,\cos\dfrac{k\pi x}{l}$ на $(0,l)$ --- для условий Неймана;
  \item полиномы Лежандра --- в сферических задачах;
  \item функции Бесселя --- в круге и цилиндре.
\end{itemize}

\textbf{Практический смысл:} ортогональность позволяет найти коэффициенты, полнота гарантирует, что базиса хватает.

\subsection{6(c). Разложение по собственным функциям задачи Штурма--Лиувилля}

Общая задача Штурма--Лиувилля:
\[
[p(x)X'(x)]'+[\lambda \rho(x)-q(x)]X(x)=0,
\qquad 0<x<l,
\]
с граничными условиями
\[
\alpha X(0)-\beta X'(0)=0,
\qquad
\gamma X(l)+\delta X'(l)=0.
\]

Если $\lambda_k\ne\lambda_m$, то собственные функции ортогональны с весом $\rho$:
\[
\int_0^l \rho(x)X_k(x)X_m(x)\,dx=0.
\]

Разложение функции имеет вид
\[
f(x)\sim \sum_{k=1}^{\infty} c_k X_k(x),
\]
где
\[
c_k=
\frac{\displaystyle\int_0^l \rho(x)f(x)X_k(x)\,dx}
{\displaystyle\int_0^l \rho(x)X_k^2(x)\,dx}.
\]

\textbf{Это обобщённый ряд Фурье:} базис задаётся не заранее, а самой краевой задачей.

\subsection{6(d). Чем разложение в ряд Фурье отличается от разложения по специальным функциям}

\textbf{Ряд Фурье} использует фиксированный базис синусов и косинусов.

\textbf{Разложение по специальным функциям} использует базис, который возникает после разделения переменных:
\begin{itemize}
  \item Лежандр --- в сферической геометрии;
  \item Бессель --- в круге и цилиндре;
  \item собственные функции Штурма--Лиувилля --- в общем случае.
\end{itemize}

\textbf{Различия:}
\begin{itemize}
  \item другой интервал ортогональности;
  \item может появляться вес $\rho(x)$;
  \item базис зависит от оператора и граничных условий.
\end{itemize}

\subsection{6(e). Физический смысл коэффициентов Фурье в задачах теплопроводности и колебаний}

Коэффициенты Фурье показывают, \textbf{насколько сильно возбуждена каждая собственная мода}.

В теплопроводности
\[
u(x,t)=\sum_{k=1}^{\infty} a_k e^{-\omega_k^2 t}X_k(x)
\]
коэффициент $a_k$ измеряет вклад $k$-й моды в начальное распределение температуры.

В колебаниях
\[
u(x,t)=\sum_{k=1}^{\infty}
\bigl(a_k\cos\omega_k t+b_k\sin\omega_k t\bigr)X_k(x)
\]
\[
a_k \leftrightarrow \text{начальное смещение},
\qquad
b_k \leftrightarrow \text{начальная скорость}.
\]

\textbf{Физический вывод:} коэффициенты говорят, какие формы колебаний или тепловые гармоники реально присутствуют в системе.

\section{Практические задачи к разделу 6}

\subsection{Задача 6(a). Разложить \texorpdfstring{$f(x)=x$}{f(x)=x} на \texorpdfstring{$(0,\pi)$}{(0,pi)} в ряд Фурье по синусам}

Ищем
\[
x\sim \sum_{k=1}^{\infty} b_k\sin kx,
\qquad
b_k=\frac{2}{\pi}\int_0^\pi x\sin kx\,dx.
\]

Интегрирование по частям даёт
\[
\int_0^\pi x\sin kx\,dx=\frac{\pi(-1)^{k+1}}{k}.
\]
Следовательно,
\[
b_k=\frac{2(-1)^{k+1}}{k}.
\]

\textbf{Ответ:}
\[
x=2\sum_{k=1}^{\infty}\frac{(-1)^{k+1}}{k}\sin kx,
\qquad 0<x<\pi.
\]

\subsection{Задача 6(b). Чем отличается разложение по синусам от разложения по косинусам и когда какое используется}

\textbf{Синусный ряд:}
\[
f(x)\sim \sum_{k=1}^{\infty}b_k\sin\frac{k\pi x}{l}.
\]
Естествен для:
\begin{itemize}
  \item нечётного продолжения;
  \item нулевых условий Дирихле.
\end{itemize}

\textbf{Косинусный ряд:}
\[
f(x)\sim \frac{a_0}{2}+\sum_{k=1}^{\infty}a_k\cos\frac{k\pi x}{l}.
\]
Естествен для:
\begin{itemize}
  \item чётного продолжения;
  \item нулевых условий Неймана.
\end{itemize}

\textbf{Короткое правило выбора:} Дирихле $\Rightarrow$ синусы, Нейман $\Rightarrow$ косинусы.

\subsection{Задача 6(c). Задача Штурма--Лиувилля для \texorpdfstring{$d^2/dx^2$}{d2/dx2} с условиями Дирихле}

Рассматриваем
\[
X''+\lambda X=0,
\qquad
X(0)=0,
\qquad
X(l)=0.
\]

\textbf{Нетривиальные решения существуют только при}
\[
\lambda_k=\left(\frac{k\pi}{l}\right)^2,
\qquad
k=1,2,\dots
\]

\textbf{Собственные функции:}
\[
X_k(x)=\sin\frac{k\pi x}{l}.
\]

Они ортогональны:
\[
\int_0^l \sin\frac{k\pi x}{l}\sin\frac{m\pi x}{l}\,dx=0
\qquad (k\ne m).
\]

\textbf{Ответ:} это основной синусный базис для метода Фурье при условиях Дирихле.

\section{Сравнение трёх типов уравнений}

\subsection{7(a). Основные свойства решений эллиптических, параболических и гиперболических уравнений}

\textbf{Эллиптические} уравнения ($\Delta u=0$, $-\Delta u=f$):
\begin{itemize}
  \item описывают стационарные поля;
  \item имеют принцип максимума;
  \item реальные характеристики в смысле переноса информации отсутствуют;
  \item решение внутри области обычно очень гладкое.
\end{itemize}

\textbf{Параболические} уравнения ($u_t-a^2\Delta u=f$):
\begin{itemize}
  \item описывают диффузию и теплопроводность;
  \item обладают сглаживающим эффектом;
  \item имеют принцип максимума в пространственно-временной области;
  \item возмущение распространяется мгновенно по всей области.
\end{itemize}

\textbf{Гиперболические} уравнения ($u_{tt}-a^2\Delta u=f$):
\begin{itemize}
  \item описывают волны и инерционные процессы;
  \item имеют конечную скорость распространения;
  \item естественно связаны с характеристиками;
  \item не обладают тем мгновенным сглаживанием, что есть у параболических задач.
\end{itemize}

\textbf{Правило-памятка:}
\[
\text{эллиптика} \leftrightarrow \text{стационарность},\qquad
\text{параболика} \leftrightarrow \text{диффузия},\qquad
\text{гиперболика} \leftrightarrow \text{волны}.
\]

\subsection{7(b). Какие типы задач ставятся для каждого типа}

\textbf{Для эллиптических уравнений:}
\begin{itemize}
  \item задача Дирихле;
  \item задача Неймана;
  \item третья краевая задача;
  \item смешанные краевые задачи.
\end{itemize}

\textbf{Для параболических уравнений:}
\begin{itemize}
  \item начально-краевые задачи;
  \item задача Коши на всей прямой или во всём пространстве.
\end{itemize}

\textbf{Для гиперболических уравнений:}
\begin{itemize}
  \item задача Коши;
  \item смешанные начально-краевые задачи;
  \item задачи на характеристиках.
\end{itemize}

\textbf{По числу начальных условий:}
\begin{itemize}
  \item эллиптический тип: задаются граничные данные;
  \item параболический тип: одно начальное условие;
  \item гиперболический тип: два начальных условия.
\end{itemize}

\end{document}
